\chapter{Powerful Questions\\คำถามที่ขยายความสามารถในการคิด}

คำถามทรงพลังไม่ได้มีคุณค่าเพราะฟังซับซ้อน แต่เพราะช่วยให้ Mentee เห็นสิ่งที่ยังไม่เห็น ทดสอบสมมติฐาน และเลือกการกระทำด้วยตนเอง คำถามที่ดีจึงไม่ใช่คำแนะนำที่เติมเครื่องหมายคำถามท้ายประโยค

\learningobjectives{
  \item ออกแบบคำถามเปิดที่มีเป้าหมายชัด
  \item เลือกลำดับคำถามจากการสำรวจสู่การตัดสินใจ
  \item หลีกเลี่ยงคำถามนำ คำถามซ้อน และบรรยากาศสอบสวน
}

\section{คุณสมบัติของคำถามที่ดี}
คำถามที่ดีสั้น เปิด และมีเพียงหนึ่งประเด็น ใช้ภาษาที่ Mentee เข้าใจ และเกิดจากสิ่งที่เพิ่งได้ยิน ไม่ใช่จากชุดคำถามที่ Mentor ต้องการถามให้ครบ คำว่า \enquote{อะไร} และ \enquote{อย่างไร} มักเปิดการสำรวจได้ดีกว่า \enquote{ทำไม} ซึ่งบางบริบทอาจทำให้รู้สึกต้องแก้ตัว

\section{คลังคำถามตามหน้าที่}
\begin{longtable}{P{.25\textwidth}P{.62\textwidth}}
\toprule
\textbf{หน้าที่} & \textbf{ตัวอย่างคำถาม}\\
\midrule
ทำให้ชัด & \enquote{เมื่อท่านพูดว่าสำเร็จ ภาพที่สังเกตได้คืออะไร}\\
ขยายมุมมอง & \enquote{ใครอีกที่มองปัญหานี้ต่างออกไป}\\
ตรวจสมมติฐาน & \enquote{ท่านกำลังสมมติอะไร และมีข้อมูลใดสนับสนุน}\\
เข้าถึงประเด็นจริง & \enquote{หากตัดเรื่องเร่งด่วนออก ประเด็นสำคัญที่สุดคืออะไร}\\
สร้างทางเลือก & \enquote{มีวิธีใดอีก แม้ยังไม่สมบูรณ์}\\
สร้างความรับผิดชอบ & \enquote{ก้าวแรกที่ท่านจะทำ เมื่อใด และเราจะรู้ได้อย่างไร}\\
\bottomrule
\end{longtable}

\section{จังหวะและลำดับ}
เริ่มจากคำถามกว้างเพื่อให้ Mentee กำหนดเรื่อง จากนั้นเจาะรายละเอียดและสมมติฐาน ก่อนเปิดทางเลือกและสรุปการกระทำ การยิงคำถามต่อเนื่องโดยไม่สะท้อนสิ่งที่ได้ยินทำให้บทสนทนากลายเป็นการสัมภาษณ์ Mentor ควรยอมให้คำถามสำคัญหนึ่งข้อมีเวลาทำงาน

\begin{examplebox}
\textbf{Mentee:} \enquote{ฉันต้องหาพันธมิตรอุตสาหกรรม}\\
\textbf{Mentor:} \enquote{พันธมิตรที่เหมาะควรเติมสิ่งใดที่ทีมยังไม่มี}\\
\textbf{Mentee:} \enquote{ช่องทางเข้าถึงผู้ใช้และความรู้เรื่องมาตรฐาน}\\
\textbf{Mentor:} \enquote{สองสิ่งนี้จำเป็นต้องมาจากองค์กรเดียวกันหรือไม่}\\
คำถามที่สองทำให้ Mentee เห็นสมมติฐานที่จำกัดทางเลือกโดยไม่จำเป็น
\end{examplebox}

\section{คำถามท้าทายอย่างเคารพ}
เมื่อจำเป็นต้องท้าทาย ให้ขออนุญาตและอิงข้อมูล เช่น \enquote{ผมขอลองเสนอข้อสังเกตที่อาจไม่สบายใจได้ไหม จากสามโครงการที่ผ่านมา ดูเหมือนทีมจะเริ่มสร้างก่อนคุยกับผู้ใช้ ท่านเห็นรูปแบบเดียวกันหรือไม่} การท้าทายเช่นนี้ตรงแต่ไม่ตัดสินตัวบุคคล

\begin{mentornote}
หากท่านรู้คำตอบที่อยากให้ Mentee พูด คำถามนั้นอาจเป็นคำถามนำ ลองเปิดเผยความเห็นของท่านอย่างตรงไปตรงมา แล้วคืนสิทธิด้วยคำถามว่า \enquote{ท่านเห็นต่างอย่างไร}
\end{mentornote}

\begin{keytakeaway}
Powerful Questions ขยายอำนาจในการคิดของ Mentee ด้วยความชัด ความอยากรู้ และจังหวะที่เหมาะ เป้าหมายไม่ใช่ถามให้มาก แต่ถามสิ่งที่ช่วยให้เกิดความเข้าใจและการเลือก
\end{keytakeaway}

\begin{reflection}
นำคำถามที่ท่านมักใช้สามข้อมาเขียนใหม่ให้สั้น เปิด และไม่ชี้คำตอบ จากนั้นกำหนดว่าคำถามแต่ละข้อมีหน้าที่ทำให้ชัด ตรวจสมมติฐาน สร้างทางเลือก หรือสร้างความรับผิดชอบ
\blankline\blankline
\end{reflection}

\chaptersummary{คำถามที่มีพลังเกิดจากการฟังและมีหน้าที่ชัด Mentor ควรเรียงคำถามจากการทำความเข้าใจไปสู่ทางเลือกและการลงมือทำ พร้อมท้าทายสมมติฐานโดยไม่ลดทอนศักดิ์ศรีของ Mentee}
